\documentclass[11pt,a4paper]{article}

% ============================================================
% FICHE D'EXERCICES - EPREUVE ANTICIPEE DE MATHEMATIQUES
% Classe de premiere - Tronc commun / voie generale hors specialite
% ============================================================

\usepackage[utf8]{inputenc}
\usepackage[T1]{fontenc}
\usepackage[french]{babel}
\usepackage{lmodern}
\usepackage{geometry}
\geometry{margin=1.7cm}
\usepackage{amsmath,amssymb,amsthm}
\usepackage{array,tabularx,booktabs,multicol}
\usepackage{enumitem}
\usepackage{xcolor}
\usepackage{tikz}
\usepackage[most]{tcolorbox}
\usepackage{fancyhdr}
\usepackage{hyperref}

\hypersetup{colorlinks=true,linkcolor=blue!60!black,urlcolor=blue!60!black}
\setlength{\parindent}{0pt}
\setlength{\parskip}{0.35em}
\setlength{\headheight}{14pt}
\renewcommand{\arraystretch}{1.25}
\setlist[enumerate]{itemsep=0.35em,topsep=0.35em}

\pagestyle{fancy}
\fancyhf{}
\lhead{Premiere -- Epreuve anticipee de mathematiques}
\rhead{Fiche d'exercices progressive}
\cfoot{\thepage}

\definecolor{bleu}{RGB}{27,75,140}
\definecolor{vert}{RGB}{20,120,75}
\definecolor{orange}{RGB}{215,110,20}
\definecolor{rouge}{RGB}{170,45,45}
\definecolor{gris}{RGB}{245,247,250}

\newtcolorbox{objectifbox}{enhanced,breakable,colback=bleu!5,colframe=bleu,title=Objectifs de la fiche,fonttitle=\bfseries}
\newtcolorbox{methodebox}{enhanced,breakable,colback=vert!5,colframe=vert,title=Methode,fonttitle=\bfseries}
\newtcolorbox{attentionbox}{enhanced,breakable,colback=rouge!5,colframe=rouge,title=Erreur classique,fonttitle=\bfseries}
\newtcolorbox{exo}[2][]{enhanced,breakable,colback=white,colframe=bleu!75!black,title={Exercice #2\ifx&#1&\else\ -- #1\fi},fonttitle=\bfseries}
\newtcolorbox{corrige}[1]{enhanced,breakable,colback=gris,colframe=black!45,title={Corrige de l'exercice #1},fonttitle=\bfseries}
\newcommand{\etoile}{\ensuremath{\star}}
\newcommand{\deuxetoiles}{\ensuremath{\star\star}}
\newcommand{\troisetoiles}{\ensuremath{\star\star\star}}
\newcommand{\quatretoiles}{\ensuremath{\star\star\star\star}}
\newcommand{\R}{\mathbb{R}}
\newcommand{\N}{\mathbb{N}}
\newcommand{\Pcond}[2]{P_{#1}\left(#2\right)}

\begin{document}

\begin{center}
{\Large\bfseries Fiche d'exercices complete et progressive}\\[0.2cm]
{\LARGE\bfseries Preparation a l'epreuve anticipee de mathematiques}\\[0.2cm]
{\large Classe de premiere -- tronc commun, voie generale hors specialite mathematiques}\\[0.2cm]
\end{center}

\begin{objectifbox}
Cette fiche est construite avec un gradient de difficulte. Elle vise a entrainer les eleves a la structure de l'epreuve :
\begin{itemize}
    \item une premiere partie d'automatismes, sans calculatrice et souvent sous forme de QCM ;
    \item une deuxieme partie d'exercices contextualises, ou il faut justifier, modeliser, calculer et interpreter ;
    \item des problemes de synthese melant pourcentages, fonctions, suites, statistiques, probabilites et derivation.
\end{itemize}
Les exercices sont classes ainsi : \etoile{} application directe, \deuxetoiles{} methode classique, \troisetoiles{} probleme type epreuve, \quatretoiles{} probleme de synthese difficile.
\end{objectifbox}

\tableofcontents
\newpage

% ============================================================
\section{Partie A -- Automatismes sans calculatrice}

\begin{methodebox}
Pour la partie automatismes, on cherche d'abord la methode la plus courte. Un bon reflexe consiste a transformer les pourcentages en coefficients multiplicateurs, a simplifier les fractions avant de calculer, a verifier les ordres de grandeur et a reconnaitre rapidement les formes usuelles.
\end{methodebox}

\subsection{QCM d'entrainement}
Pour chaque question, une seule reponse est correcte. Aucune justification n'est exigee dans une vraie partie QCM, mais la correction donne la methode.

\begin{exo}[Automatismes de base \etoile]{A1}
Repondre aux questions suivantes.
\begin{enumerate}
    \item Un prix de 80 euros augmente de 25\%. Le nouveau prix est :\quad a) 85 \quad b) 100 \quad c) 105 \quad d) 120.
    \item Diminuer une quantite de 12\% revient a la multiplier par :\quad a) 1,12 \quad b) 0,12 \quad c) 0,88 \quad d) 1,88.
    \item Un prix augmente de 10\% puis diminue de 10\%. Le prix final est :\quad a) identique \quad b) plus grand \quad c) plus petit \quad d) impossible a savoir.
    \item Le taux d'evolution reciproque d'une hausse de 25\% est environ :\quad a) $-25\%$ \quad b) $-20\%$ \quad c) $+20\%$ \quad d) $-30\%$.
    \item $\dfrac{3}{4}+\dfrac{5}{8}$ vaut :\quad a) $\dfrac{8}{12}$ \quad b) $\dfrac{11}{8}$ \quad c) $\dfrac{13}{8}$ \quad d) $\dfrac{1}{8}$.
    \item $2^{-3}$ vaut :\quad a) $-8$ \quad b) $\dfrac18$ \quad c) $\dfrac12$ \quad d) $8$.
    \item $10^3\times 10^{-5}$ vaut :\quad a) $10^{-15}$ \quad b) $10^{-2}$ \quad c) $10^8$ \quad d) $10^2$.
    \item Si $4$ objets coutent $18$ euros, alors $10$ objets coutent :\quad a) 36 \quad b) 40 \quad c) 45 \quad d) 50 euros.
    \item $35\%$ de $200$ vaut :\quad a) 7 \quad b) 35 \quad c) 70 \quad d) 700.
    \item Une vitesse de $90$ km/h correspond a :\quad a) $15$ m/s \quad b) $20$ m/s \quad c) $25$ m/s \quad d) $30$ m/s.
\end{enumerate}
\end{exo}

\begin{exo}[Automatismes fonctions et statistiques \etoile]{A2}
Repondre aux questions suivantes.
\begin{enumerate}
    \item La droite passant par $A(1;3)$ et $B(5;11)$ a pour coefficient directeur : a) 1 b) 2 c) 3 d) 4.
    \item La droite d'equation $y=-2x+5$ coupe l'axe des ordonnees au point : a) $(0;5)$ b) $(5;0)$ c) $(0;-2)$ d) $(-2;0)$.
    \item Pour $f(x)=3x^2-2x+1$, l'image de $-1$ est : a) 0 b) 2 c) 6 d) 10.
    \item La forme developpee de $(x-4)(x+3)$ est : a) $x^2-x-12$ b) $x^2+x-12$ c) $x^2-7x+12$ d) $x^2-12$.
    \item Une racine de $x^2-5x+6$ est : a) 1 b) 2 c) 5 d) 6.
    \item La moyenne de $4,6,10$ est : a) 6 b) $\dfrac{20}{3}$ c) 7 d) 8.
    \item La mediane de $2,2,8,10,15$ est : a) 2 b) 8 c) 10 d) 15.
    \item Dans une classe de 30 eleves, 18 sont des filles. La proportion de filles est : a) 0,3 b) 0,4 c) 0,6 d) 0,8.
    \item Si $P(A)=0,7$, alors $P(\overline A)$ vaut : a) 0,7 b) 0,3 c) 1,7 d) impossible.
    \item Si $P(A)=0,4$ et $P_A(B)=0,25$, alors $P(A\cap B)$ vaut : a) 0,65 b) 0,15 c) 0,10 d) 0,25.
\end{enumerate}
\end{exo}

\begin{exo}[Calcul mental dirige \deuxetoiles]{A3}
Calculer rapidement.
\begin{multicols}{2}
\begin{enumerate}
    \item $101\times 99$.
    \item $1,2^2$.
    \item $0,8^3$.
    \item $\dfrac{15}{100}\times 240$.
    \item $\dfrac{2}{3}\div \dfrac{5}{6}$.
    \item $\sqrt[3]{1000}$.
    \item $\sqrt[4]{81}$.
    \item $3(2x-5)-2(x+1)$.
    \item $(x+5)^2-(x-5)^2$.
    \item $\dfrac{3V}{\pi r^2}=h$ : isoler $V$.
\end{enumerate}
\end{multicols}
\end{exo}

% ============================================================
\section{Partie B -- Pourcentages, proportions et evolutions}

\begin{methodebox}
Pour une evolution de taux $t$, le coefficient multiplicateur est $1+t$ si $t$ est ecrit sous forme decimale. Une hausse de $15\%$ correspond a $1,15$. Une baisse de $15\%$ correspond a $0,85$. Pour plusieurs evolutions successives, on multiplie les coefficients multiplicateurs.
\end{methodebox}

\begin{attentionbox}
Une hausse de $20\%$ suivie d'une baisse de $20\%$ ne ramene pas a la valeur initiale. En effet, $1,20\times 0,80=0,96$ : la valeur finale est inferieure de $4\%$ a la valeur initiale.
\end{attentionbox}

\begin{exo}[Coefficients multiplicateurs \etoile]{B1}
Completer le tableau.
\begin{center}
\begin{tabular}{|c|c|c|}
\hline
Evolution & Coefficient multiplicateur & Exemple pour une valeur initiale de 200 \\
\hline
Hausse de $8\%$ & & \\
\hline
Baisse de $8\%$ & & \\
\hline
Hausse de $35\%$ & & \\
\hline
Baisse de $2,5\%$ & & \\
\hline
\end{tabular}
\end{center}
\end{exo}

\begin{exo}[Evolution successive \deuxetoiles]{B2}
Un abonnement coute 50 euros. Il augmente de $12\%$ en janvier, puis diminue de $5\%$ en juillet.
\begin{enumerate}
    \item Calculer le prix final.
    \item Calculer le coefficient multiplicateur global.
    \item En deduire le taux d'evolution global.
\end{enumerate}
\end{exo}

\begin{exo}[Taux reciproque \deuxetoiles]{B3}
Une population augmente de $40\%$.
\begin{enumerate}
    \item Par quel coefficient a-t-elle ete multipliee ?
    \item Quel taux de baisse faudrait-il appliquer ensuite pour revenir a la valeur initiale ?
    \item Expliquer pourquoi ce taux n'est pas $40\%$.
\end{enumerate}
\end{exo}

\begin{exo}[Tableau d'evolution \troisetoiles]{B4}
Une entreprise vend 12 000 produits en 2024. Les ventes diminuent de $10\%$ en 2025 puis augmentent de $25\%$ en 2026.
\begin{enumerate}
    \item Calculer le nombre de produits vendus en 2025 puis en 2026.
    \item Calculer le taux d'evolution global entre 2024 et 2026.
    \item Une eleve affirme : « Comme $-10+25=15$, les ventes ont augmente de $15\%$. » Corriger son raisonnement.
\end{enumerate}
\end{exo}

% ============================================================
\section{Partie C -- Fonctions affines, droites et modeles lineaires}

\begin{methodebox}
Une fonction affine est de la forme $f(x)=mx+p$. Le nombre $m$ est le coefficient directeur : il mesure la variation de $f(x)$ lorsque $x$ augmente d'une unite. Le nombre $p=f(0)$ est l'ordonnee a l'origine. Pour deux points $A(x_A;y_A)$ et $B(x_B;y_B)$ avec $x_A\neq x_B$, on a
\[
 m=\frac{y_B-y_A}{x_B-x_A}.
\]
\end{methodebox}

\begin{exo}[Lire et utiliser une droite \etoile]{C1}
On considere la fonction affine $f(x)=2,5x-4$.
\begin{enumerate}
    \item Calculer $f(0)$, $f(2)$ et $f(10)$.
    \item Resoudre $f(x)=0$.
    \item Interpreter le coefficient directeur.
\end{enumerate}
\end{exo}

\begin{exo}[Determiner une fonction affine \deuxetoiles]{C2}
Une droite passe par les points $A(2;7)$ et $B(6;19)$.
\begin{enumerate}
    \item Calculer son coefficient directeur.
    \item Determiner son equation reduite.
    \item Le point $C(10;31)$ appartient-il a cette droite ?
\end{enumerate}
\end{exo}

\begin{exo}[Modele lineaire de refroidissement \troisetoiles]{C3}
Un liquide est a $90^\circ$C a l'instant $t=0$. Au bout de 4 minutes, il est a $70^\circ$C. On suppose, dans un premier modele, que la temperature diminue de facon lineaire.
\begin{enumerate}
    \item Determiner une fonction affine $T(t)$ donnant la temperature en fonction du temps $t$ en minutes.
    \item Selon ce modele, au bout de combien de temps la temperature atteint-elle $40^\circ$C ?
    \item Pourquoi ce modele devient-il discutable pour de grandes valeurs de $t$ ?
\end{enumerate}
\end{exo}

% ============================================================
\section{Partie D -- Statistiques, tableaux croises et information chiffree}

\begin{methodebox}
Dans un tableau croise, on distingue les effectifs totaux, les effectifs marginaux, les frequences globales et les frequences conditionnelles. Dire « parmi les filles, 60\% ont choisi anglais » ne signifie pas la meme chose que « parmi les eleves ayant choisi anglais, 60\% sont des filles ».
\end{methodebox}

\begin{exo}[Indicateurs statistiques \etoile]{D1}
On considere la serie suivante :
\[
8,\ 9,\ 10,\ 10,\ 11,\ 12,\ 20.
\]
\begin{enumerate}
    \item Calculer la moyenne.
    \item Determiner la mediane.
    \item Determiner l'etendue.
    \item Expliquer pourquoi la moyenne et la mediane donnent ici deux informations differentes.
\end{enumerate}
\end{exo}

\begin{exo}[Tableau croise \deuxetoiles]{D2}
Dans un lycee de 500 eleves, on etudie le choix d'une option.
\begin{center}
\begin{tabular}{|c|c|c|c|}
\hline
 & Option musique & Pas musique & Total \\
\hline
Filles & 80 & 170 & 250 \\
\hline
Garcons & 50 & 200 & 250 \\
\hline
Total & 130 & 370 & 500 \\
\hline
\end{tabular}
\end{center}
\begin{enumerate}
    \item Quelle est la proportion d'eleves ayant choisi musique ?
    \item Quelle est la proportion de filles parmi les eleves ayant choisi musique ?
    \item Quelle est la proportion d'eleves ayant choisi musique parmi les filles ?
    \item Ces deux proportions precedent-elles la meme information ?
\end{enumerate}
\end{exo}

\begin{exo}[Boite a moustaches et comparaison \deuxetoiles]{D3}
Deux classes obtiennent les notes suivantes a un controle.
\[
A: 7,8,9,10,10,11,12,13,14,16
\]
\[
B: 3,5,8,10,10,10,12,15,18,19
\]
\begin{enumerate}
    \item Calculer la moyenne de chaque classe.
    \item Determiner la mediane de chaque classe.
    \item Comparer la dispersion des deux series a l'aide de l'etendue.
    \item Quelle classe semble la plus homogene ?
\end{enumerate}
\end{exo}

\begin{exo}[Filtres logiques \troisetoiles]{D4}
On dispose d'une base de donnees d'eleves avec trois informations : sexe, option, presence a une sortie. On note :
\[
F : \text{« l'eleve est une fille »},\quad M : \text{« l'eleve a choisi musique »},\quad S : \text{« l'eleve participe a la sortie »}.
\]
Traduire en langage courant :
\begin{enumerate}
    \item $F\cap M$.
    \item $M\cup S$.
    \item $F\cap \overline{S}$.
    \item $\overline{M\cup S}$.
\end{enumerate}
Puis traduire en notation ensembliste :
\begin{enumerate}[resume]
    \item « L'eleve est une fille et ne participe pas a la sortie. »
    \item « L'eleve ne fait ni musique ni sortie. »
\end{enumerate}
\end{exo}

% ============================================================
\section{Partie E -- Probabilites conditionnelles et independance}

\begin{methodebox}
Si $P(A)\neq 0$, la probabilite de $B$ sachant $A$ est
\[
P_A(B)=\frac{P(A\cap B)}{P(A)}.
\]
Deux evenements $A$ et $B$ sont independants lorsque $P(A\cap B)=P(A)P(B)$, ou encore lorsque $P_A(B)=P(B)$ si $P(A)\neq 0$.
\end{methodebox}

\begin{exo}[Arbre pondere \deuxetoiles]{E1}
Dans une association, $60\%$ des adherents sont adultes. Parmi les adultes, $30\%$ pratiquent la natation. Parmi les mineurs, $50\%$ pratiquent la natation. On choisit un adherent au hasard.
\begin{enumerate}
    \item Representer la situation par un arbre pondere.
    \item Calculer la probabilite que l'adherent soit adulte et pratique la natation.
    \item Calculer la probabilite que l'adherent pratique la natation.
    \item Sachant que l'adherent pratique la natation, calculer la probabilite qu'il soit mineur.
\end{enumerate}
\end{exo}

\begin{exo}[Independance \deuxetoiles]{E2}
On donne $P(A)=0,4$, $P(B)=0,5$ et $P(A\cap B)=0,2$.
\begin{enumerate}
    \item Calculer $P_A(B)$.
    \item Les evenements $A$ et $B$ sont-ils independants ?
\end{enumerate}
\end{exo}

\begin{exo}[Lancers independants \troisetoiles]{E3}
Une piece truquee donne pile avec probabilite $0,3$. On la lance trois fois de suite, les lancers etant independants.
\begin{enumerate}
    \item Calculer la probabilite d'obtenir exactement trois piles.
    \item Calculer la probabilite de n'obtenir aucun pile.
    \item Calculer la probabilite d'obtenir exactement une pile.
    \item Calculer la probabilite d'obtenir au moins une pile.
\end{enumerate}
\end{exo}

% ============================================================
\section{Partie F -- Suites et modeles discrets}

\begin{methodebox}
Une suite arithmetique verifie $u_{n+1}=u_n+r$, donc $u_n=u_0+nr$. Elle modelise une evolution a accroissement constant. Une suite geometrique verifie $u_{n+1}=q u_n$, donc $u_n=u_0q^n$. Elle modelise une evolution a taux constant.
\end{methodebox}

\begin{exo}[Suite arithmetique \etoile]{F1}
Une salle contient 120 chaises. Chaque semaine, on ajoute 15 chaises.
\begin{enumerate}
    \item Definir une suite $u_n$ modelisant le nombre de chaises apres $n$ semaines.
    \item Donner la nature de la suite, son premier terme et sa raison.
    \item Exprimer $u_n$ en fonction de $n$.
    \item Au bout de combien de semaines y aura-t-il au moins 300 chaises ?
\end{enumerate}
\end{exo}

\begin{exo}[Suite geometrique \deuxetoiles]{F2}
Une population de bacteries est de 500 au depart et augmente de $18\%$ par heure.
\begin{enumerate}
    \item Definir une suite $v_n$ modelisant la population apres $n$ heures.
    \item Donner sa nature et sa raison.
    \item Exprimer $v_n$ en fonction de $n$.
    \item Calculer $v_5$ a l'unite pres.
\end{enumerate}
\end{exo}

\begin{exo}[Comparer deux modeles \troisetoiles]{F3}
Un lac contient initialement $1000$ plantes invasives. Deux modeles sont proposes.
\begin{itemize}
    \item Modele A : le nombre de plantes augmente de $300$ par mois.
    \item Modele B : le nombre de plantes augmente de $20\%$ par mois.
\end{itemize}
\begin{enumerate}
    \item Donner l'expression du nombre de plantes apres $n$ mois dans chaque modele.
    \item Calculer les valeurs pour $n=0,1,2,3,6$.
    \item A partir de quel mois le modele B depasse-t-il le modele A ?
    \item Expliquer pourquoi une croissance geometrique peut depasser une croissance arithmetique meme si elle parait plus lente au debut.
\end{enumerate}
\end{exo}

\begin{exo}[Seuil et algorithme \troisetoiles]{F4}
Une suite est definie par $u_0=250$ et $u_{n+1}=1,08u_n$.
\begin{enumerate}
    \item Que modelise une telle recurrence ?
    \item Ecrire un algorithme en langage naturel qui determine le plus petit entier $n$ tel que $u_n\geq 1000$.
    \item Determiner ce plus petit entier a l'aide de calculs successifs.
\end{enumerate}
\end{exo}

% ============================================================
\section{Partie G -- Croissance exponentielle et racines n-iemes}

\begin{methodebox}
Lorsque $a>0$, la fonction $x\mapsto a^x$ prolonge l'idee des suites geometriques. Si une grandeur est multipliee par $a$ a chaque unite de temps, alors on la modelise par $f(t)=C a^t$. La racine $n$-ieme d'un nombre positif $A$ est le nombre positif $r$ tel que $r^n=A$.
\end{methodebox}

\begin{exo}[Puissances et racines \etoile]{G1}
Calculer.
\begin{multicols}{2}
\begin{enumerate}
    \item $2^5$.
    \item $3^4$.
    \item $10^{-2}$.
    \item $\sqrt[3]{8}$.
    \item $\sqrt[4]{16}$.
    \item $\sqrt[5]{32}$.
\end{enumerate}
\end{multicols}
\end{exo}

\begin{exo}[Taux annuel moyen \deuxetoiles]{G2}
Un capital passe de 1000 euros a 1331 euros en 3 ans, avec un taux annuel constant.
\begin{enumerate}
    \item On note $q$ le coefficient multiplicateur annuel. Justifier que $q^3=1,331$.
    \item Determiner $q$.
    \item En deduire le taux annuel moyen.
\end{enumerate}
\end{exo}

\begin{exo}[Modele exponentiel continu simple \troisetoiles]{G3}
On modelise une quantite par $f(t)=80\times 1,15^t$, ou $t$ est exprime en annees.
\begin{enumerate}
    \item Calculer $f(0)$, $f(1)$ et $f(2)$.
    \item Interpreter $1,15$.
    \item Determiner le taux d'evolution entre $t=0$ et $t=2$.
    \item Trouver un coefficient multiplicateur semestriel $q$ tel que deux semestres donnent une hausse annuelle de $15\%$.
\end{enumerate}
\end{exo}

% ============================================================
\section{Partie H -- Derivation de polynomes et variations}

\begin{methodebox}
Pour un polynome de degre inferieur ou egal a 3, on utilise les regles :
\[
(ax^3+bx^2+cx+d)'=3ax^2+2bx+c.
\]
Le signe de $f'$ donne les variations de $f$ : si $f'>0$, alors $f$ est croissante ; si $f'<0$, alors $f$ est decroissante.
\end{methodebox}

\begin{exo}[Derivees directes \etoile]{H1}
Calculer les derivees des fonctions suivantes.
\begin{enumerate}
    \item $f(x)=4x^2-3x+7$.
    \item $g(x)=-2x^3+5x^2-x+1$.
    \item $h(x)=\dfrac12x^3-4x$.
\end{enumerate}
\end{exo}

\begin{exo}[Variations d'un trinome \deuxetoiles]{H2}
On considere $f(x)=x^2-6x+5$.
\begin{enumerate}
    \item Calculer $f'(x)$.
    \item Etudier le signe de $f'(x)$.
    \item Dresser le tableau de variations de $f$.
    \item En deduire la valeur minimale de $f$.
\end{enumerate}
\end{exo}

\begin{exo}[Optimisation \troisetoiles]{H3}
On veut fabriquer une boite ouverte a partir d'une feuille carree de cote 20 cm. On decoupe dans chaque coin un carre de cote $x$, puis on replie les bords.
\begin{enumerate}
    \item Expliquer pourquoi $0<x<10$.
    \item Montrer que le volume de la boite est $V(x)=x(20-2x)^2$.
    \item Developper $V(x)$.
    \item Calculer $V'(x)$.
    \item Determiner les variations de $V$ sur $]0;10[$.
    \item En deduire la valeur de $x$ donnant le volume maximal.
\end{enumerate}
\end{exo}

% ============================================================
\section{Partie I -- Problemes type epreuve}

\begin{exo}[Probleme 1 : refroidissement d'un plat \troisetoiles]{I1}
Un plat sort du four a $180^\circ$C. La temperature de la piece est $25^\circ$C. On note $T_n$ la temperature du plat $n$ minutes apres sa sortie du four.

\textbf{Modele 1.} On suppose que la temperature diminue de $25^\circ$C par minute.
\begin{enumerate}
    \item Exprimer $T_n$ en fonction de $n$.
    \item Selon ce modele, calculer $T_5$ et $T_8$.
    \item Critiquer ce modele.
\end{enumerate}

\textbf{Modele 2.} On note $U_n=T_n-25$ l'ecart entre la temperature du plat et celle de la piece. On suppose que cet ecart diminue de $18\%$ par minute.
\begin{enumerate}[resume]
    \item Donner $U_0$.
    \item Justifier que $U_{n+1}=0,82U_n$.
    \item Exprimer $U_n$ puis $T_n$ en fonction de $n$.
    \item Le plat peut etre servi lorsque $T_n\leq 45$. Determiner a partir de quelle minute il peut etre servi.
\end{enumerate}
\end{exo}

\begin{exo}[Probleme 2 : sport et probabilites \troisetoiles]{I2}
Lors d'une journee sportive, $70\%$ des participants choisissent la course, les autres choisissent la marche. Parmi les coureurs, $40\%$ sont licencies. Parmi les marcheurs, $10\%$ sont licencies. On choisit une personne au hasard.

On note $C$ : « la personne choisit la course » et $L$ : « la personne est licenciee ».
\begin{enumerate}
    \item Donner $P(C)$, $P(\overline C)$, $P_C(L)$ et $P_{\overline C}(L)$.
    \item Construire un arbre pondere.
    \item Calculer $P(C\cap L)$ et $P(\overline C\cap L)$.
    \item Calculer $P(L)$.
    \item Sachant que la personne est licenciee, calculer la probabilite qu'elle ait choisi la course.
\end{enumerate}
\end{exo}

\begin{exo}[Probleme 3 : bacteries et seuil sanitaire \quatretoiles]{I3}
La concentration de bacteries dans une piscine est initialement de $800$ bacteries par millilitre. Le seuil maximal autorise est de $1500$ bacteries par millilitre.

\textbf{Modele A.} On suppose que la concentration augmente de $90$ bacteries par millilitre chaque heure.
\begin{enumerate}
    \item Exprimer la concentration $A_n$ apres $n$ heures.
    \item Au bout de combien d'heures le seuil est-il depasse ?
\end{enumerate}

\textbf{Modele B.} On suppose que la concentration augmente de $12\%$ par heure.
\begin{enumerate}[resume]
    \item Exprimer la concentration $B_n$ apres $n$ heures.
    \item Calculer $B_4$ et $B_5$ a l'unite pres.
    \item Au bout de combien d'heures le seuil est-il depasse ?
    \item Quel modele est le plus inquietant a long terme ? Justifier.
\end{enumerate}

\textbf{Modele C.} On propose la fonction $f(t)=800+100t+5t^2$ pour modeliser la concentration apres $t$ heures.
\begin{enumerate}[resume]
    \item Calculer $f'(t)$.
    \item Etudier les variations de $f$ sur $[0;6]$.
    \item Resoudre par essais encadres l'inequation $f(t)>1500$ sur $[0;6]$.
\end{enumerate}
\end{exo}

\begin{exo}[Probleme 4 : donnees, probabilites, independance \quatretoiles]{I4}
Dans un lycee de 1200 eleves, on donne le tableau suivant.
\begin{center}
\begin{tabular}{|c|c|c|c|}
\hline
 & Pratique un sport & Ne pratique pas de sport & Total \\
\hline
Filles & 300 & 260 & 560 \\
\hline
Garcons & 400 & 240 & 640 \\
\hline
Total & 700 & 500 & 1200 \\
\hline
\end{tabular}
\end{center}
On choisit un eleve au hasard. On note $F$ l'evenement « l'eleve est une fille » et $S$ l'evenement « l'eleve pratique un sport ».
\begin{enumerate}
    \item Calculer $P(F)$, $P(S)$ et $P(F\cap S)$.
    \item Calculer $P_F(S)$ et $P_S(F)$.
    \item Les evenements $F$ et $S$ sont-ils independants ? Justifier de deux manieres.
    \item Traduire en francais $F\cap \overline S$ et calculer sa probabilite.
    \item Une personne affirme : « Comme il y a plus de sportifs garcons que de sportives filles, un sportif a plus d'une chance sur deux d'etre un garcon. » Cette affirmation est-elle correcte ? Justifier.
\end{enumerate}
\end{exo}

% ============================================================
\section{Partie J -- Devoir bilan final}

\begin{exo}[Devoir bilan -- Partie 1 : automatismes, 6 points]{J1}
Pour chaque question, choisir la bonne reponse.
\begin{enumerate}
    \item $25\%$ de 360 vaut : a) 45 b) 60 c) 90 d) 120.
    \item Une baisse de $7\%$ correspond au coefficient : a) 1,07 b) 0,07 c) 0,93 d) 0,97.
    \item $\dfrac35+\dfrac{1}{10}$ vaut : a) $\dfrac{4}{15}$ b) $\dfrac{7}{10}$ c) $\dfrac{2}{5}$ d) $\dfrac{6}{15}$.
    \item Le coefficient directeur de la droite passant par $(2;5)$ et $(6;13)$ vaut : a) 1 b) 2 c) 3 d) 4.
    \item Si $f(x)=-x^2+4x+1$, alors $f(2)$ vaut : a) 1 b) 4 c) 5 d) 9.
    \item Une quantite multipliee par $1,21$ a augmente de : a) $12,1\%$ b) $21\%$ c) $121\%$ d) $0,21\%$.
    \item $\sqrt[3]{64}$ vaut : a) 4 b) 6 c) 8 d) 16.
    \item Si $P(A)=0,35$, alors $P(\overline A)$ vaut : a) 0,35 b) 0,65 c) 1,35 d) impossible.
    \item La mediane de $4,4,5,9,20$ est : a) 4 b) 5 c) 9 d) 20.
    \item $10^{-3}$ vaut : a) $-1000$ b) $0,001$ c) $0,003$ d) $1000$.
    \item Si $u_{n+1}=1,04u_n$, la suite modelise : a) une baisse de 4\% b) une hausse de 4\% c) une hausse de 104\% d) une evolution lineaire.
    \item La derivee de $3x^2-5x+1$ est : a) $3x-5$ b) $6x-5$ c) $6x+1$ d) $x^3-5x$.
\end{enumerate}
\end{exo}

\begin{exo}[Devoir bilan -- Partie 2 : probleme complet, 14 points]{J2}
Une commune suit la surface occupee par une plante invasive dans un etang de $5000$ m$^2$. Au debut de l'etude, la plante occupe $400$ m$^2$.

\textbf{Partie A -- Modele lineaire.} On suppose que la surface augmente de $180$ m$^2$ par semaine.
\begin{enumerate}
    \item Definir une suite $u_n$ donnant la surface apres $n$ semaines.
    \item Exprimer $u_n$ en fonction de $n$.
    \item Au bout de combien de semaines la surface depasse-t-elle $2000$ m$^2$ ?
\end{enumerate}

\textbf{Partie B -- Modele geometrique.} On suppose maintenant que la surface augmente de $18\%$ par semaine.
\begin{enumerate}[resume]
    \item Definir une suite $v_n$ donnant la surface apres $n$ semaines.
    \item Exprimer $v_n$ en fonction de $n$.
    \item Calculer $v_5$ a l'unite pres.
    \item A l'aide d'essais successifs, determiner au bout de combien de semaines la surface depasse $2000$ m$^2$.
\end{enumerate}

\textbf{Partie C -- Intervention.} La commune choisit au hasard une parcelle de l'etang. On sait que $42\%$ des parcelles sont proches de la rive. Parmi les parcelles proches de la rive, $30\%$ sont envahies. Parmi les autres parcelles, $12\%$ sont envahies.

On note $R$ : « la parcelle est proche de la rive » et $E$ : « la parcelle est envahie ».
\begin{enumerate}[resume]
    \item Donner les probabilites directement lues dans l'enonce.
    \item Construire un arbre pondere.
    \item Calculer $P(R\cap E)$.
    \item Calculer $P(E)$.
    \item Sachant qu'une parcelle est envahie, calculer la probabilite qu'elle soit proche de la rive.
\end{enumerate}

\textbf{Partie D -- Fonction d'intervention.} Le cout d'une intervention, en milliers d'euros, est modelise par
\[
C(x)=0,02x^3-0,6x^2+7x+15,
\]
ou $x$ designe le nombre d'hectares traites, avec $0\leq x\leq 10$.
\begin{enumerate}[resume]
    \item Calculer $C'(x)$.
    \item Montrer que $C'(x)=0,06x^2-1,2x+7$ est toujours positif sur $[0;10]$ en calculant son minimum.
    \item En deduire les variations de $C$ sur $[0;10]$.
\end{enumerate}
\end{exo}

\newpage
% ============================================================
\section{Corrections detaillees}

\begin{corrige}{A1}
\begin{enumerate}
    \item $80\times 1,25=100$. Reponse b.
    \item Une baisse de $12\%$ donne $1-0,12=0,88$. Reponse c.
    \item $1,10\times0,90=0,99$, donc le prix final est plus petit. Reponse c.
    \item Apres une hausse de $25\%$, on multiplie par $1,25$. Pour revenir au depart, il faut multiplier par $1/1,25=0,8$, soit une baisse de $20\%$. Reponse b.
    \item $\frac34=\frac68$, donc $\frac34+\frac58=\frac{11}{8}$. Reponse b.
    \item $2^{-3}=1/2^3=1/8$. Reponse b.
    \item $10^3\times10^{-5}=10^{-2}$. Reponse b.
    \item Un objet coute $18/4=4,5$ euros, donc $10$ coutent $45$ euros. Reponse c.
    \item $35\%$ de 200 vaut $0,35\times200=70$. Reponse c.
    \item $90$ km/h $=90\times1000/3600=25$ m/s. Reponse c.
\end{enumerate}
\end{corrige}

\begin{corrige}{A2}
\begin{enumerate}
    \item $m=(11-3)/(5-1)=8/4=2$. Reponse b.
    \item Pour $x=0$, $y=5$. Reponse a.
    \item $f(-1)=3+2+1=6$. Reponse c.
    \item $(x-4)(x+3)=x^2-x-12$. Reponse a.
    \item $2^2-5\times2+6=0$. Reponse b.
    \item Moyenne : $(4+6+10)/3=20/3$. Reponse b.
    \item La valeur centrale est $8$. Reponse b.
    \item $18/30=0,6$. Reponse c.
    \item $P(\overline A)=1-P(A)=0,3$. Reponse b.
    \item $P(A\cap B)=P(A)P_A(B)=0,4\times0,25=0,10$. Reponse c.
\end{enumerate}
\end{corrige}

\begin{corrige}{A3}
\begin{multicols}{2}
\begin{enumerate}
    \item $101\times99=(100+1)(100-1)=10000-1=9999$.
    \item $1,2^2=1,44$.
    \item $0,8^3=0,512$.
    \item $15\%$ de $240$ vaut $36$.
    \item $\frac23\div\frac56=\frac23\times\frac65=\frac45$.
    \item $\sqrt[3]{1000}=10$.
    \item $\sqrt[4]{81}=3$.
    \item $3(2x-5)-2(x+1)=6x-15-2x-2=4x-17$.
    \item Difference de carres : $(x+5)^2-(x-5)^2=20x$.
    \item $V=\dfrac{\pi r^2h}{3}$? Ici $\frac{3V}{\pi r^2}=h$, donc $3V=h\pi r^2$, puis $V=\frac{h\pi r^2}{3}$.
\end{enumerate}
\end{multicols}
\end{corrige}

\begin{corrige}{B1}
Hausse de $8\%$ : coefficient $1,08$, valeur $216$. Baisse de $8\%$ : coefficient $0,92$, valeur $184$. Hausse de $35\%$ : coefficient $1,35$, valeur $270$. Baisse de $2,5\%$ : coefficient $0,975$, valeur $195$.
\end{corrige}

\begin{corrige}{B2}
\begin{enumerate}
    \item $50\times1,12\times0,95=53,20$. Le prix final est $53,20$ euros.
    \item Coefficient global : $1,12\times0,95=1,064$.
    \item Le taux global est $1,064-1=0,064$, soit $6,4\%$.
\end{enumerate}
\end{corrige}

\begin{corrige}{B3}
\begin{enumerate}
    \item Hausse de $40\%$ : coefficient $1,4$.
    \item Pour revenir au depart, il faut multiplier par $1/1,4=5/7\approx0,7143$. Le taux est donc $-28,57\%$ environ.
    \item $40\%$ de la valeur augmentee est plus grand que $40\%$ de la valeur initiale. Les pourcentages ne s'additionnent pas directement.
\end{enumerate}
\end{corrige}

\begin{corrige}{B4}
\begin{enumerate}
    \item $2025:12000\times0,90=10800$. $2026:10800\times1,25=13500$.
    \item Coefficient global $0,90\times1,25=1,125$, donc hausse globale de $12,5\%$.
    \item L'erreur consiste a additionner les taux. Il faut multiplier les coefficients multiplicateurs.
\end{enumerate}
\end{corrige}

\begin{corrige}{C1}
\begin{enumerate}
    \item $f(0)=-4$, $f(2)=1$, $f(10)=21$.
    \item $2,5x-4=0$, donc $2,5x=4$ et $x=1,6$.
    \item Quand $x$ augmente de $1$, $f(x)$ augmente de $2,5$.
\end{enumerate}
\end{corrige}

\begin{corrige}{C2}
\begin{enumerate}
    \item $m=(19-7)/(6-2)=12/4=3$.
    \item $y=3x+p$. Avec $A(2;7)$ : $7=6+p$, donc $p=1$. Equation : $y=3x+1$.
    \item Pour $x=10$, $3x+1=31$. Donc $C$ appartient a la droite.
\end{enumerate}
\end{corrige}

\begin{corrige}{C3}
\begin{enumerate}
    \item Le coefficient directeur est $(70-90)/(4-0)=-5$. Donc $T(t)=90-5t$.
    \item $90-5t=40$, donc $t=10$ minutes.
    \item Pour de grandes valeurs, le modele donne des temperatures inferieures a celle de la piece, ce qui n'est pas realiste pour un refroidissement simple.
\end{enumerate}
\end{corrige}

\begin{corrige}{D1}
Somme : $80$, effectif : $7$, moyenne $80/7\approx11,43$. La mediane est la 4e valeur, donc $10$. L'etendue vaut $20-8=12$. La valeur $20$ tire la moyenne vers le haut, alors que la mediane reste representative du centre de la serie.
\end{corrige}

\begin{corrige}{D2}
\begin{enumerate}
    \item $130/500=0,26$, soit $26\%$.
    \item Parmi les eleves ayant choisi musique, la proportion de filles vaut $80/130=8/13$.
    \item Parmi les filles, la proportion ayant choisi musique vaut $80/250=0,32$.
    \item Non. La premiere conditionne par « musique », la seconde par « filles ».
\end{enumerate}
\end{corrige}

\begin{corrige}{D3}
La moyenne de $A$ vaut $110/10=11$. Celle de $B$ vaut $110/10=11$. Les deux medianes valent $(10+11)/2=10,5$ pour $A$ et $(10+10)/2=10$ pour $B$. L'etendue de $A$ vaut $16-7=9$, celle de $B$ vaut $19-3=16$. La classe $A$ est plus homogene car son etendue est plus faible.
\end{corrige}

\begin{corrige}{D4}
\begin{enumerate}
    \item $F\cap M$ : l'eleve est une fille et a choisi musique.
    \item $M\cup S$ : l'eleve a choisi musique ou participe a la sortie, eventuellement les deux.
    \item $F\cap\overline S$ : l'eleve est une fille et ne participe pas a la sortie.
    \item $\overline{M\cup S}$ : l'eleve ne fait ni musique ni sortie.
    \item $F\cap\overline S$.
    \item $\overline M\cap\overline S$, ou encore $\overline{M\cup S}$.
\end{enumerate}
\end{corrige}

\begin{corrige}{E1}
\begin{enumerate}
    \item Arbre : $A$ avec $0,60$, $\overline A$ avec $0,40$ ; depuis $A$, natation $0,30$ ; depuis $\overline A$, natation $0,50$.
    \item $0,60\times0,30=0,18$.
    \item $P(N)=0,60\times0,30+0,40\times0,50=0,18+0,20=0,38$.
    \item $P_N(\overline A)=0,20/0,38=10/19\approx0,526$.
\end{enumerate}
\end{corrige}

\begin{corrige}{E2}
$P_A(B)=P(A\cap B)/P(A)=0,2/0,4=0,5$. Comme $P_A(B)=P(B)=0,5$, les evenements sont independants. On peut aussi verifier $P(A)P(B)=0,4\times0,5=0,2=P(A\cap B)$.
\end{corrige}

\begin{corrige}{E3}
\begin{enumerate}
    \item $0,3^3=0,027$.
    \item $0,7^3=0,343$.
    \item Exactement une pile : trois positions possibles, donc $3\times0,3\times0,7^2=0,441$.
    \item Au moins une pile : $1-P(aucune)=1-0,343=0,657$.
\end{enumerate}
\end{corrige}

\begin{corrige}{F1}
$u_0=120$ et $u_{n+1}=u_n+15$. C'est une suite arithmetique de raison $15$. Donc $u_n=120+15n$. On cherche $120+15n\geq300$, soit $15n\geq180$, donc $n\geq12$. Il faut 12 semaines.
\end{corrige}

\begin{corrige}{F2}
$v_0=500$ et $v_{n+1}=1,18v_n$. C'est une suite geometrique de raison $1,18$. Donc $v_n=500\times1,18^n$. On obtient $v_5\approx500\times2,2878\approx1144$ bacteries.
\end{corrige}

\begin{corrige}{F3}
Modele A : $A_n=1000+300n$. Modele B : $B_n=1000\times1,2^n$. Valeurs :
\[
\begin{array}{c|ccccc}
n&0&1&2&3&6\\ \hline
A_n&1000&1300&1600&1900&2800\\
B_n&1000&1200&1440&1728&2986
\end{array}
\]
A $n=6$, le modele B depasse le modele A. La croissance geometrique augmente d'autant plus vite que la valeur deja atteinte est grande.
\end{corrige}

\begin{corrige}{F4}
La recurrence modelise une hausse de $8\%$ a chaque etape. Algorithme : initialiser $n=0$, $u=250$ ; tant que $u<1000$, remplacer $u$ par $1,08u$ et $n$ par $n+1$ ; afficher $n$. Calculs successifs : $u_n=250\times1,08^n$. On cherche $1,08^n\geq4$. Par essais : $1,08^{18}\approx3,996$, $1,08^{19}\approx4,316$. Le plus petit entier est $19$.
\end{corrige}

\begin{corrige}{G1}
$2^5=32$, $3^4=81$, $10^{-2}=0,01$, $\sqrt[3]{8}=2$, $\sqrt[4]{16}=2$, $\sqrt[5]{32}=2$.
\end{corrige}

\begin{corrige}{G2}
On a $1000q^3=1331$, donc $q^3=1,331$. Or $1,1^3=1,331$, donc $q=1,1$. Le taux annuel moyen est $10\%$.
\end{corrige}

\begin{corrige}{G3}
$f(0)=80$, $f(1)=92$, $f(2)=105,8$. Le nombre $1,15$ signifie une hausse de $15\%$ par an. Entre $0$ et $2$, le coefficient est $1,15^2=1,3225$, soit une hausse de $32,25\%$. Le coefficient semestriel $q$ verifie $q^2=1,15$, donc $q=\sqrt{1,15}\approx1,072$.
\end{corrige}

\begin{corrige}{H1}
$f'(x)=8x-3$, $g'(x)=-6x^2+10x-1$, $h'(x)=\frac32x^2-4$.
\end{corrige}

\begin{corrige}{H2}
$f'(x)=2x-6=2(x-3)$. Donc $f'<0$ si $x<3$, $f'=0$ si $x=3$, $f'>0$ si $x>3$. La fonction decroit sur $]-\infty;3]$ puis croit sur $[3;+\infty[$. Minimum : $f(3)=9-18+5=-4$.
\end{corrige}

\begin{corrige}{H3}
On doit avoir $0<x<10$ pour que la hauteur soit positive et que la base ait un cote positif. La base mesure $20-2x$, donc $V(x)=x(20-2x)^2=x(400-80x+4x^2)=4x^3-80x^2+400x$. Ainsi $V'(x)=12x^2-160x+400=4(3x^2-40x+100)$. Les zeros de $3x^2-40x+100$ sont $x=\frac{40\pm20}{6}$, soit $x=\frac{10}{3}$ et $x=10$. Sur $]0;10[$, $V'$ est positive sur $]0;10/3[$ puis negative sur $]10/3;10[$. Le volume maximal est obtenu pour $x=10/3$ cm.
\end{corrige}

\begin{corrige}{I1}
Modele 1 : $T_n=180-25n$. $T_5=55$ et $T_8=-20$, ce qui est absurde car le plat ne peut pas descendre sous la temperature de la piece dans ce modele simple. Modele 2 : $U_0=180-25=155$. Une baisse de $18\%$ donne un coefficient $0,82$, donc $U_{n+1}=0,82U_n$ et $U_n=155\times0,82^n$. Ainsi $T_n=25+155\times0,82^n$. On cherche $T_n\leq45$, soit $155\times0,82^n\leq20$. Par essais : $n=10$ donne environ $46,3$ degres, $n=11$ donne environ $42,5$ degres. Le plat peut etre servi a partir de la 11e minute.
\end{corrige}

\begin{corrige}{I2}
$P(C)=0,70$, $P(\overline C)=0,30$, $P_C(L)=0,40$, $P_{\overline C}(L)=0,10$. Donc $P(C\cap L)=0,70\times0,40=0,28$ et $P(\overline C\cap L)=0,30\times0,10=0,03$. Ainsi $P(L)=0,31$. Sachant $L$, la probabilite d'avoir choisi la course vaut $P_L(C)=0,28/0,31=28/31\approx0,903$.
\end{corrige}

\begin{corrige}{I3}
Modele A : $A_n=800+90n$. On cherche $800+90n>1500$, soit $90n>700$, donc $n>7,77$. Le seuil est depasse au bout de 8 heures. Modele B : $B_n=800\times1,12^n$. $B_4\approx1259$ et $B_5\approx1410$ ; $B_6\approx1579$, donc seuil depasse au bout de 6 heures. A long terme, le modele geometrique est plus inquietant car la hausse est proportionnelle a la quantite presente. Modele C : $f'(t)=100+10t$, positif sur $[0;6]$, donc $f$ est croissante. $f(5)=800+500+125=1425$ et $f(6)=800+600+180=1580$, donc le depassement a lieu entre 5 h et 6 h.
\end{corrige}

\begin{corrige}{I4}
$P(F)=560/1200=7/15$, $P(S)=700/1200=7/12$, $P(F\cap S)=300/1200=1/4$. $P_F(S)=300/560=15/28$, tandis que $P_S(F)=300/700=3/7$. Les evenements ne sont pas independants car $P(F)P(S)=\frac{7}{15}\times\frac{7}{12}=\frac{49}{180}\neq\frac14$. On peut aussi dire que $P_F(S)\neq P(S)$. $F\cap\overline S$ signifie « l'eleve est une fille et ne pratique pas de sport » ; sa probabilite vaut $260/1200=13/60$. Parmi les sportifs, il y a $400$ garcons sur $700$, soit $4/7>1/2$ : l'affirmation est correcte.
\end{corrige}

\begin{corrige}{J1}
Reponses : 1c, 2c, 3b, 4b, 5c, 6b, 7a, 8b, 9b, 10b, 11b, 12b.
\end{corrige}

\begin{corrige}{J2}
Partie A. $u_0=400$ et $u_{n+1}=u_n+180$, donc $u_n=400+180n$. On cherche $400+180n>2000$, soit $180n>1600$, donc $n>8,88$. Il faut 9 semaines.

Partie B. $v_0=400$ et $v_{n+1}=1,18v_n$, donc $v_n=400\times1,18^n$. $v_5\approx400\times2,2878\approx915$. Par essais : $v_9\approx1772$, $v_{10}\approx2091$. Le seuil de $2000$ m$^2$ est depasse au bout de 10 semaines.

Partie C. $P(R)=0,42$, $P(\overline R)=0,58$, $P_R(E)=0,30$, $P_{\overline R}(E)=0,12$. $P(R\cap E)=0,42\times0,30=0,126$. $P(E)=0,126+0,58\times0,12=0,126+0,0696=0,1956$. Donc $P_E(R)=0,126/0,1956\approx0,644$.

Partie D. $C'(x)=0,06x^2-1,2x+7$. Cette parabole a son sommet en $x=-b/(2a)=1,2/(0,12)=10$. Son minimum sur $[0;10]$ vaut donc $C'(10)=6-12+7=1>0$. Ainsi $C'(x)>0$ sur $[0;10]$, donc $C$ est strictement croissante sur $[0;10]$.
\end{corrige}

\newpage
\section{Fiche de synthese des reflexes}

\begin{methodebox}
\textbf{Pourcentages.} Remplacer les taux par des coefficients multiplicateurs. Plusieurs evolutions successives se traitent par multiplication, jamais par addition directe des taux.

\textbf{Fonctions affines.} Penser a $f(x)=mx+p$ et a $m=\dfrac{\Delta y}{\Delta x}$. Un modele affine traduit une variation absolue constante.

\textbf{Statistiques.} Distinguer moyenne, mediane et etendue. Distinguer proportion globale et proportion conditionnelle.

\textbf{Probabilites.} Sur un arbre, on multiplie le long d'une branche et on additionne les branches correspondant au meme evenement. Pour une probabilite conditionnelle, identifier clairement l'univers restreint.

\textbf{Suites.} Une augmentation d'une meme quantite donne une suite arithmetique. Une augmentation d'un meme pourcentage donne une suite geometrique.

\textbf{Exponentielle.} $a^x$ prolonge la logique geometrique. La racine $n$-ieme sert a retrouver un coefficient moyen.

\textbf{Derivation.} La derivee mesure une variation instantanee. Son signe donne les variations de la fonction.
\end{methodebox}

\end{document}
